\scanpage{61}
% manuscript line(s) 1–4
d'un élément d'un torseur sous $L^{\otimes2}\otimes\check E$ \inserted{$=L^{\otimes2}\otimes\mathcal A$}, qui \deleted{\unsure{est}} peut être considéré comme l'espace des applications \uline{affines} de $P_1$ dans $L^{\otimes2}$, \inserted{i.e. des applications linéaires de $E$ dans $L^{\otimes2}$, soit $\check E\otimes L^{\otimes2}$} \deleted{\unsure{donc c'est}} et on a
% manuscript line(s) 5
\[
0\unsure{\mathbin{-}|\to}L^{\otimes2}\longrightarrow L^{\otimes2}\otimes\check E\mathbin{-}L\longrightarrow0
\]
% manuscript line(s) 5–7
\margin{(car $0\mathbin{-}\underline O_X\mathbin{-}\check E\mathbin{-}\check L\longrightarrow0$) et on multiplie par $L^{\otimes2}$}
% manuscript line(s) 6–8
$\underset{L\otimes E}{\unsure{\shortparallel|?|\simeq}}$ \uline{can} (pour tt module loc. lib. de rg $2$, \unsure{$F$|$E$} avec $L\simeq\det E$) i.e.
% manuscript line(s) 8
\[
E\simeq\check E\otimes L\ \big(\simeq\underline{\mathrm{Hom}}(E,\underset{\substack{\shortparallel\\\det E}}L)\big)
\]

% manuscript line(s) 9
L'application affine
% manuscript line(s) 10
\[
P_1\longrightarrow L^{\otimes2}
\]
% manuscript line(s) 11
définie par un $\unsure{\zeta|\xi}\in\underset{\substack{\cap\\E}}L\cdot E\simeq L\otimes E\subset\mathrm{Sym}^2E$
% manuscript line(s) 12
est définie par
% manuscript line(s) 13
\[
x\longmapsto x\deleted{\unsure{\otimes}}\wedge\zeta
\]
% manuscript line(s) 12–15
\margin{$E\otimes E\to L$ \quad $(x\otimes y)\mapsto x\wedge y$ \quad $\underbrace{L\otimes E}_{\ni\,\zeta}\otimes\underset{\substack{\ni\\x}}E\mathbin{-}L\otimes L=L^{\otimes2}$ \quad $(\zeta,x)\mathbin{-}\zeta\wedge x=-x\wedge\zeta$}
% manuscript line(s) 16
$\ldots\,.$]

% manuscript line(s) 17–18
Inversement, si (l'ext. $E$ de $\underline O_X$ par $L$ étant donnée, sans $Q$ encore) si on se donne
% manuscript line(s) 19
\[
T\in E,\ N\in\mathrm{Sym}^2(E),
\]
% manuscript line(s) 20–21
quand correspondent-ils à une $Q$ comme dans 3) ou 3$'$) ? Il faut, rappelons le
% manuscript line(s) 22
\[
\psi_1(T)=2,\qquad\unsure{\psi_2|\varphi_2}(N)=1
\]
% manuscript line(s) 22–23
\margin{i.e. (en termes de la ``section unité'' $e$ de $\mathcal A=\check E$ : $T(e)=2,N(e)=1$)}
% manuscript line(s) 24
plus une condition qui équivaudrait à
% manuscript line(s) 25
\[
N(x+\deleted{\unsure{\lambda}}e)=N(x)+\deleted{\unsure{\lambda}}T(x)+1^{\deleted{\unsure{2}}}
\]
% manuscript line(s) 25
\margin{(\inserted{on aura $N(x+\lambda e)=N(x)+\lambda T(x)+\lambda^2$})}
% manuscript line(s) 26
i.e. $T(x)=\varphi_N(x,e)$

% manuscript line(s) 27
En tous cas on doit avoir : \uline{$N$ détermine $T$}
% archive sheet number: 60
